\section{源码节点全覆盖索引}

Lecture 6 是 executable source，真正的“slide-complete”不是逐页 PDF，而是逐个教学函数覆盖。下面这张表说明每个函数在讲义中的位置和处理方式。

\begin{longtable}{p{0.20\textwidth}p{0.30\textwidth}p{0.40\textwidth}}
\toprule
\textbf{源码函数} & \textbf{教学目标} & \textbf{讲义处理} \\
\midrule
\texttt{review\_of\_gpus} & 复习硬件与编程模型 & GPU 图、grid 图、block occupancy 图、术语表。 \\
\texttt{benchmarking} & 端到端计时 & CUDA events、warmup、synchronize 代码。 \\
\texttt{profiling} & 看 kernel 级时间分布 & profiler 表解释、kernel name 解读。 \\
GeLU variants & fusion 的性能收益 & GeLU 公式、naive/builtin/compiled 对比。 \\
Triton GeLU & elementwise kernel & block offsets、mask、load/store、PTX 观察。 \\
Triton softmax & row-wise reduction & max-subtract-exp-sum-divide 的 fused row kernel。 \\
Triton row sum & row too large for one block & tile loop、accumulator、final reduction。 \\
Triton matmul ReLU & tiled GEMM + fusion & tile pointers、\texttt{tl.dot}、fused ReLU。 \\
\bottomrule
\end{longtable}

\begin{importantbox}{这张索引表的目的}
它是本讲的 coverage matrix 的正文版本。读者可以看到：本讲不是泛泛讲 Triton，而是沿源码中的每个教学函数建立 benchmark → profile → optimize → write kernels 的完整链条。
\end{importantbox}

